\documentclass[10pt,a4paper]{amsart}
\usepackage[margin=19mm]{geometry}
\usepackage{amsmath,amssymb,mathtools}
\usepackage[scr=rsfs,cal=euler]{mathalfa}
\usepackage{xfrac}
\usepackage[unicode,hidelinks,bookmarks=false]{hyperref}
\newcommand{\ie}{i.e.\ }
\newcommand{\cf}{cf.\ }
\newcommand{\resp}{respectively\ }
\newcommand{\Subsection}{\subsection}
\providecommand{\nobreakdash}{\nobreak\hbox{-}}
\setlength{\emergencystretch}{2em}
\multlinegap0pt
\usepackage{amssymb}
\usepackage{bm}
\usepackage{tikz}
\usepackage{float}
\usepackage{enumerate}
\usepackage{xcolor}
\usepackage{mathtools}
\newtheorem{theorem}{Theorem}
\title{Graham conjecture on small sets in abelian groups}
\author{Simone Costa, Stefano Della Fiore, Mattia Fontana, Llu\'is Vena}
\date{Source: arXiv:2603.20961v1 (CC BY 4.0)}
\begin{document}
\maketitle

\noindent\emph{Source and licence.} Graham conjecture on small sets in abelian groups. Version arXiv:2603.20961v1 (CC BY 4.0), distributed by arXiv under the Creative Commons Attribution 4.0 International licence, http://creativecommons.org/licenses/by/4.0/, as recorded in the arXiv licence metadata for this exact version and shown on the abstract page https://arxiv.org/abs/2603.20961v1. Full licence notice: \texttt{licenses/arxiv-2603.20961-CC-BY-4.0.txt}.\par\smallskip

\noindent\textit{Editorial scope.} Counted unit: the article's theorem prop:another\_3 (source0.tex lines 349--351). The article's complete original proof of it is delivered with it (source0.tex lines 353--392, one \texttt{proof} environment of 40 lines). No mathematics is written, repaired or re-derived: the statement and the proof are the article's own lines, in the article's own notation, and no retained line is substituted or re-flowed.  No further source-local context is needed: every cross-reference of every retained line resolves inside the two delivered spans.  Nothing was omitted from any retained span, so the package carries no omission record.  No retained line contains a citation, so the wrapper carries no bibliography and no bibliography entry.  No retained line contains a figure environment, an image-inclusion command or drawing code, so no image is delivered and the package declares no asset dependency.  Editorial formatting only, in addition: an AMS \texttt{amsart} preamble that restates the article's own declarations and macros used by the retained lines, namely the environments \texttt{theorem} on separate counters, together with the standard packages those lines need.  The retained environments print in this document as Theorem 1 in order of appearance, whereas the article prints the same items with the numbering of its own full document.  The complete native source of the article as distributed in the arXiv e-print for this exact version is bundled unchanged as \texttt{sources/196899/source0.tex}.
\begin{theorem}[Characterization Theorem]\label{prop:another_3}
	Let $G$ be an abelian group and let $A=\{a_1,\ldots,a_k\}$ $\subseteq G\setminus \{0\}$ be a set of $k\geq 2$ elements such that, for any distinct $i,j$, $a_i+a_j\in \{0,a_1,\ldots,a_k\}$. Then $A\cup \{0\}$  is a subgroup of $G$.
\end{theorem}

\begin{proof}[Proof of Theorem~\ref{prop:another_3}]
First of all, we consider the case where $A$ is constituted only by involutions. In this case, it is clear that $A\cup \{0\}$ is inverse-closed and it is closed under the sum of its elements and is hence a subgroup of $G$. Similarly, also when $|A|=2$ we have that $A\cup \{0\}$ is inverse-closed and it is closed under the sum. So in the following we will assume the existence of a non-involution element and that $|A|\geq 3$.

Now we prove that, under these hypotheses, there exists a non-involution $a_i\in A$ so that $-a_i\in A$.
%Otherwise, we would have that, for any distinct $\{i,j\}$, $a_i+a_j\not=0$.  This implies that, for any distinct $i,j$, we get $$a_i+a_j\in \{a_1,\ldots,a_k\}.$$ There are ${k \choose 2}$ such pairs and only $n$ elements in $A$, so each element of $A$ must be the sum of an average of  $$\frac{{k \choose 2}}{k}=(k-1)/2$$ distinct pairs. Since any $a_i$ can not be the sum of more than this number of pairs, each element of $A$ must be the sum of exactly $(k-1)/2$ distinct pairs. Note that this is not possible if $k$ is even, and thus in this case we must have a pair of distinct $\{i,j\}$ such that $a_i+a_j=0$.
%what if k is odd?
Indeed, assume otherwise. With the hypothesis for any distinct $i,j$, $a_i+a_j\in \{0,a_1,\ldots,a_k\}$ we conclude that for any distinct $\{i,j\}$, $a_i+a_j\not=0$. There are ${k \choose 2}$ such pairs and only $n$ elements in $A$, so each element of $A$ must be the sum of an average of  $$\frac{{k \choose 2}}{k}=(k-1)/2$$ distinct pairs.
If a sum would have more than this number of pairs, then two such pairs would have an element in common, thus making the two pairs the same. So each element of $A$ must be the sum of exactly $(k-1)/2$ distinct pairs. Note that this is not possible if $k$ is even. For the case $k$ odd, for each $i$ there is a matching $\{(r_i,s_i)\}_{i\in[(k-1)/2]}$ between elements in $\{a_1,\ldots,a_k\}\setminus a_i$ such that $a_{r_j}+a_{s_j}=a_i$. By considering the relation of the non-involution element $a_t$ and its matched element $a_{t'}$ and $a_i$ so that (so $a_t+a_{t'}=a_i$ with $(t,t')=(r_j,s_j)$ for some $j\in [(k-1)/2]$), and the relation $a_{t'}=a_i+a_{i'}$ for the $i'$-th element matching the $i$-th, we conclude that $a_t+a_{i'}=0$. Since $a_t$ was a non-involution element, this proves that there exists a non-involution $a_t\in A$ so that $-a_t\in A$. 


%Also, given a non-involution element $a_i$, we may assume (up to renaming the elements) that $a_i\not=a_2$ and $a_1 = a_2+a_i$.  Moreover, there exists also an index $l$ such that $$a_2 = a_1 + a_l.$$ Substituting $a_1 = a_2 + a_i$ into this equation yields $$a_2 = (a_2 + a_i) + a_l,$$ hence $a_i + a_l = 0$. This proves that there exists a non-involution $a_i\in A$ so that $-a_i\in A$. 

Now, assuming that $a_i,-a_i\in A$ with $a_i\neq -a_i$, we prove that $2a_j\in A\cup \{0\}$ for all $a_j\not=\pm a_i$. Indeed, both $a_j+a_i$ and $a_j-a_i\in A\cup \{0\}$ and, since these elements are distinct, also
$$(a_j+a_i)+(a_j-a_i)=2a_j\in A\cup \{0\},$$
and claim follows.
Now, since $2a_j=a_j+a_j\in A\cup \{0\}$, it follows that, for all $a_j\not=\pm a_i$, the sets 
$$\{0+a_j,a_1+a_j,\dots,a_k+a_j\}\subseteq A\cup \{0\}.$$
Since both these sets have cardinality $k+1$, we must have that 
$$\{0+a_j,a_1+a_j,\dots,a_k+a_j\}= A\cup \{0\}$$
and hence $-a_j\in A$ for all $j$ (also when $j=i$).

Finally, we consider $2a_i$ and want to see that $2a_i\in A\cup \{0\}$. We assume $a_1\not=\pm a_i$, and we recall that we have 
$\{0+a_1,a_1+a_1,\dots,a_k+a_1\}= A\cup \{0\}$.
Therefore, there exists $j$ such that $a_i=a_1+a_j$. 
If $a_j= a_i$, then we have $a_1=0$ a contradiction, and if $a_j= -a_i$ then $a_1=2a_i$ as desired. If $a_j\not=\pm a_i$, then 
%It is clear that also $a_j\not=\pm a_i$. 
%why it is clear?
%This means that 
this means that
$$2a_i=a_i+a_i=(a_j+a_1)+a_i=a_j+(a_1+a_i).$$
Note that there exists an index $h$ such that $a_1+a_i=a_i+a_1=a_h\in A$ which implies that 
$$2a_i=a_j+(a_1+a_i)=a_j+a_h=a_h+a_j \in A\cup \{0\}$$
where the last inclusion holds since $a_j\not=\pm a_i$. We also note that the same argument (exchanging $a_i$ by $-a_i$) applies for $2(-a_i)=-2a_i\in A$.
We have proved that
$2a_j\in A\cup\{0\}$ for all $j$ (also when considering $\pm a_i$). It follows that $A\cup \{0\}$ is both inverse-closed and closed under the sums of its elements and thus is a subgroup.
%Finally, we consider $\pm 2a_i$. We assume $a_1\not=\pm a_i$, and we recall that we have  $\{0+a_1,a_1+a_1,\dots,a_k+a_1\}= A\cup \{0\}$. Therefore, there exists $j$ such that $a_i=a_1+a_j$.  If $a_j= \pm a_i$, then we have $a_1=-2a_i$ (as $a_1\neq 0$) and were are done. If $a_j\not=\pm a_i$, then it is clear that also $a_j\not=\pm a_i$. 
%why it is clear?
%This means that 
%this means that $$2a_i=a_i+a_i=(a_j+a_1)+a_i=a_j+(a_1+a_i).$$ Note that there exists an index $h$ such that $a_1+a_i=a_i+a_1=a_h\in A$ which implies that  $$2a_i=a_j+(a_1+a_i)=a_j+a_h=a_h+a_j \in A\cup \{0\}$$ where the last inclusion holds since $a_j\not=\pm a_i$. We also note that $2(-a_i)=-2a_i\in A$ since it is the additive inverse of $2a_i$. We have proved that $2a_j\in A\cup\{0\}$ for all $j$ (also when considering $\pm a_i$). It follows that $A\cup \{0\}$ is both inverse-closed and closed under the sums of its elements and thus is a subgroup.
 \end{proof}

\end{document}
